\documentclass[10pt,a4paper]{amsart}
\usepackage[margin=19mm]{geometry}
\usepackage{amsmath,amssymb,mathtools}
\usepackage[scr=rsfs,cal=euler]{mathalfa}
\usepackage{xfrac}
\usepackage[unicode,hidelinks,bookmarks=false]{hyperref}
\newcommand{\ie}{i.e.\ }
\newcommand{\cf}{cf.\ }
\newcommand{\resp}{respectively\ }
\newcommand{\Subsection}{\subsection}
\providecommand{\nobreakdash}{\nobreak\hbox{-}}
\setlength{\emergencystretch}{2em}
\multlinegap0pt
\usepackage{mathtools}
\usepackage{bm}
\usepackage[shortlabels]{enumitem}
\usepackage{xcolor}
\usepackage{tcolorbox}
\newtheorem{lemma}{Lemma}
\newtheorem{theorem}{Theorem}
\newcommand{\R}{\mathbb{R}}
\title{Persistence Pairing Graphs: Tracking Changes of Selected Homology Bases}
\author{John Rick Manzanares}
\date{Source: arXiv:2608.29719v1 (CC BY 4.0)}
\begin{document}
\maketitle

\noindent\emph{Source and licence.} Persistence Pairing Graphs: Tracking Changes of Selected Homology Bases. Version arXiv:2608.29719v1 (CC BY 4.0), distributed by arXiv under the Creative Commons Attribution 4.0 International licence, http://creativecommons.org/licenses/by/4.0/, as recorded in the arXiv licence metadata for this exact version and shown on the abstract page https://arxiv.org/abs/2608.29719v1. Full licence notice: \texttt{licenses/arxiv-2608.29719-CC-BY-4.0.txt}.\par\smallskip

\noindent\textit{Editorial scope.} Counted unit: the article's theorem thm:active-basis (source0.tex lines 737--744). The article's complete original proof of it is delivered with it (source0.tex lines 746--806, one \texttt{proof} environment of 61 lines). No mathematics is written, repaired or re-derived: the statement and the proof are the article's own lines, in the article's own notation, and no retained line is substituted or re-flowed.  Retained uncounted as the source-local context the proof needs: lines 648--655; lines 698--701.  Nothing was omitted from any retained span, so the package carries no omission record.  No retained line contains a citation, so the wrapper carries no bibliography and no bibliography entry.  No retained line contains a figure environment, an image-inclusion command or drawing code, so no image is delivered and the package declares no asset dependency.  Editorial formatting only, in addition: an AMS \texttt{amsart} preamble that restates the article's own declarations and macros used by the retained lines, namely the environments \texttt{lemma}, \texttt{theorem} on separate counters, together with the standard packages those lines need.  The retained environments print in this document as Lemma 1, Theorem 1 in order of appearance, whereas the article prints the same items with the numbering of its own full document.  The complete native source of the article as distributed in the arXiv e-print for this exact version is bundled unchanged as \texttt{sources/208882/source0.tex}.
\begin{lemma}
\label{lem:reduced-matrix-facts}
Let $U\in\Bbbk^{N\times N}$ be upper unitriangular.
\begin{enumerate}
\item If $c\in\Bbbk^N$ is nonzero, then $\operatorname{low}(Uc)=\operatorname{low}(c)$.
\item The nonzero columns of a reduced matrix are linearly independent.
\end{enumerate}
\end{lemma}

\begin{lemma}
\label{lem:pivot-positive}
If $\operatorname{low}(R_j)=i$, then $R_i=0$ and $\dim\sigma_j=\dim\sigma_i+1$.
\end{lemma}

\begin{theorem}
\label{thm:active-basis}
For every stage $p$ and every dimension $q$,
\[
\mathcal B_q(p) := \{[z_i]_p\mid i\in\mathcal A_q(p)\}
\]
is a basis of $H_q(K^p;\Bbbk)$.
\end{theorem}

\begin{proof}
Consider the columns $V_i$ satisfying $i\leq p$ and $\dim\sigma_i=q$. Because $V$ is upper unitriangular and preserves chain dimension, these columns form a basis of $C_q(K^p;\Bbbk)$. Hence, every $q$-chain $c\in C_q(K^p;\Bbbk)$ can be written uniquely as
\[
c = \sum_{\substack{i\leq p\\ \dim\sigma_i=q}}a_iV_i.
\]
Applying the boundary matrix gives
\[
Dc = \sum_{\substack{i\leq p\\ \dim\sigma_i=q}}a_iDV_i = \sum_{\substack{i\leq p\\ \dim\sigma_i=q}} a_iR_i.
\]
The nonzero columns of $R$ are linearly independent. Therefore, $Dc=0$ if and only if $a_i=0$ for every $i$ with $R_i\neq 0$. Consequently,
\[
\mathcal{Z} := \left\{z_i\ \mid \ i \leq p,\;\dim\sigma_i=q, \text{ and }R_i=0\right\}
\]
is a basis of $Z_q(K^p;\Bbbk)$.

The columns $V_j$ satisfying $j\leq p$ and $\dim\sigma_j=q+1$ form a basis of $C_{q+1}(K^p;\Bbbk)$. Their boundaries are
\[
DV_j=R_j.
\]
Hence, the nonzero columns
\[
\left\{R_j\ \mid j\leq p,\;\dim\sigma_j=q+1,\text{ and } R_j\neq 0\right\}
\]
span $B_q(K^p;\Bbbk)$. They are linearly independent because $R$ is reduced, so they form a basis of the boundary space.

Let $R_j\neq 0$ with $j\leq p$ and $\dim\sigma_j=q+1$, and set $i:=\operatorname{low}(R_j)$. Since $R_j$ is a $q$-boundary, it is also a $q$-cycle. Hence, it has a unique expansion in the cycle basis $\mathcal{Z}$.

Let $\lambda:=V^{-1}R_j$. By Lemma~\ref{lem:reduced-matrix-facts},
\[
\operatorname{low}(\lambda) = \operatorname{low}(R_j) = i.
\]
Therefore, $\lambda_i\neq 0$ and $\lambda_k=0$ for all $k>i$.

Since only positive cycle basis vectors can occur in the expansion of a cycle, we may write
\[
R_j = \lambda_i z_i + \sum_{\substack{k<i\\ R_k=0}}\lambda_k z_k,
\]
where $\lambda_i\neq 0$. By Lemma~\ref{lem:pivot-positive}, the leading index $i$ is positive.

Distinct nonzero columns $R_j$ have distinct low indices because $R$ is reduced. Thus, the boundary relations have distinct leading positive indices. These leading indices are precisely the positive indices $i$ satisfying $\operatorname{pair}(i)\leq p$.

For such an index $i$, let $j=\operatorname{pair}(i)$. Since
\[
R_j = \lambda_i z_i + \sum_{\substack{k<i\\R_k=0}} \lambda_k z_k
\]
given $\lambda_i \neq 0$ and $R_j$ is a boundary in $K^p$, its homology class is zero. Hence,
\[
0 = \lambda_i[z_i]_p + \sum_{\substack{k<i\\R_k=0}}\lambda_k[z_k]_p.
\]
Since $\lambda_i\neq0$, we obtain
\[
[z_i]_p = -\lambda_i^{-1}\sum_{\substack{k<i\\R_k=0}}\lambda_k[z_k]_p.
\]
Thus, whenever the paired death cell of $i$ has appeared, the class $[z_i]_p$ is no longer an independent homology class, that is, it is determined by classes with smaller indices.

The remaining basis classes are precisely those satisfying $i\leq p<\operatorname{pair}(i)$, which are exactly those with $i\in\mathcal A_q(p)$. Hence,
\[
\mathcal B_q(p) = \{[z_i]_p\mid i\in\mathcal A_q(p)\}
\]
is a basis of $H_q(K^p;\Bbbk)$.
\end{proof}

\end{document}
